% Part: first-order-logic
% Chapter: tableaux
% Section: provability-consistency

\documentclass[../../../include/open-logic-section]{subfiles}

\begin{document}

\iftag{FOL}
      {\olfileid{fol}{tab}{prv}}
      {\olfileid{pl}{tab}{prv}}

\olsection{\usetoken{S}{derivability} lan Konsistensi}

Saiki kita bakal netepake sawetara sipat relasi !!{derivability}.
Sipat kasebut narik kawigaten dhewe-dhewe, nanging saben sipat bakal
duwe peran ing bukti teorema kelengkapan.

\begin{prop}\ollabel{prop:provability-contr}
  Yen $\Gamma \Proves !A$ lan $\Gamma \cup \{!A\}$ ora konsisten,
  mula $\Gamma$ ora konsisten.
\end{prop}

\begin{proof}
Ana $\Gamma_0 = \{!B_1, \dots, !B_n\}$ lan $\Gamma_1
  =\{!C_1, \dots, !C_m\} \subseteq \Gamma$ kang winates saengga
  \begin{align*}
    \{\sFmla{\False}{!A}, &
    \sFmla{\True}{!B_1}, \dots, \sFmla{\True}{!B_n}\} \\
    \{\sFmla{\True}{!A}, &
    \sFmla{\True}{!C_1}, \dots, \sFmla{\True}{!C_m}\}
  \end{align*}
  nduweni !!{tableau}s katutup. Kanthi nggunakake aturan \Cut{}
  marang $!A$, kita bisa nggabungake loro tableau kasebut dadi siji
  !!{tableau} katutup kang nuduhake yen $\Gamma_0 \cup \Gamma_1$ ora
  konsisten. Amarga $\Gamma_0 \subseteq \Gamma$ lan $\Gamma_1
  \subseteq \Gamma$, kita duwe $\Gamma_0 \cup \Gamma_1 \subseteq
  \Gamma$, mula $\Gamma$ ora konsisten. Cathetan koreksi sumber
  OLPL-017: sumber Inggris menehi indeks pungkasan kang padha karo
  dhaptar kapisan marang dhaptar kapindho, nanging kabeh panganggone
  sabanjure menehi indeks pungkasan dhewe; indeks kapindho iku kang
  dibalekake ing kene.
\end{proof}

\begin{prop}
\ollabel{prop:prov-incons}
$\Gamma \Proves !A$ yen lan mung yen $\Gamma \cup \{\lnot !A\}$ ora
konsisten.
\end{prop}

\begin{proof}
Sepisan, upamana $\Gamma \Proves !A$, yaiku ana
!!{tableau} katutup kanggo
\[
\{\sFmla{\False}{!A},
\sFmla{\True}{!B_1}, \dots, \sFmla{\True}{!B_n}\}
\]
Kanthi nggunakake aturan $\TRule{\True}{\lnot}$, tableau iki bisa
dadi !!{tableau} katutup kanggo
\[
\{\sFmla{\True}{\lnot !A},
\sFmla{\True}{!B_1}, \dots, \sFmla{\True}{!B_n}\}.
\]

Kosok baline, yen ana !!{tableau} katutup kanggo himpunan kang
pungkasan mau, kita bisa ngowahi dadi !!{tableau} katutup kanggo
himpunan kapisan kanthi mbusak saben formula kang diasilake saka
\TRule{\True}{\lnot} kang ditrapake marang asumsi kapisan
$\sFmla{\True}{\lnot !A}$, uga mbusak asumsi kasebut, banjur nambahake
asumsi $\sFmla{\False}{!A}$. Sebab, yen sadurunge sawijining cabang
katutup amarga ngemot kesimpulan \TRule{\True}{\lnot} kang ditrapake
marang $\sFmla{\True}{\lnot !A}$, yaiku $\sFmla{\False}{!A}$, cabang
kang cocog ing !!{tableau} anyar uga katutup. Yen sawijining cabang
ing tableau lawas katutup amarga ngemot asumsi
$\sFmla{\True}{\lnot !A}$ uga $\sFmla{\False}{\lnot !A}$, kita bisa
ngowahi dadi cabang katutup kanthi ngetrapake
$\TRule{\False}{\lnot}$ marang $\sFmla{\False}{\lnot !A}$ kanggo
entuk $\sFmla{\True}{!A}$. Iki nutup cabang amarga kita nambahake
$\sFmla{\False}{!A}$ minangka asumsi.
\end{proof}

\begin{prob}
Buktekna yen $\Gamma \Proves \lnot !A$ yen lan mung yen $\Gamma \cup
\{!A\}$ ora konsisten.
\end{prob}

\begin{prop}\ollabel{prop:explicit-inc}
  Yen $\Gamma \Proves !A$ lan $\lnot !A \in \Gamma$, mula $\Gamma$
  ora konsisten.
\end{prop}

\begin{proof}
  Upamana $\Gamma \Proves !A$ lan $\lnot !A \in \Gamma$. Mula ana
  $!B_1$, \dots, $!B_n \in \Gamma$ saengga
  \[
  \{\sFmla{\False}{!A}, \sFmla{\True}{!B_1}, \dots, \sFmla{\True}{!B_n}\}
  \]
  nduweni tableau katutup. Gantinen asumsi
  \sFmla{\False}{!A} nganggo \sFmla{\True}{\lnot !A}, banjur
  lebokna kesimpulan saka \TRule{\True}{\lnot} kang ditrapake marang
  \sFmla{\True}{\lnot !A} sawise asumsi-asumsi. Saben !!{sentence}
  ing !!{tableau} kang kabenerake kanthi nuduhake larik~$1$ ing
  !!{tableau} lawas saiki kabenerake kanthi nuduhake larik~$n+2$.
  Dadi, yen !!{tableau} lawas katutup, tableau anyar uga katutup.
  Tableau anyar nuduhake yen $\Gamma$ ora konsisten, amarga kabeh
  asumsine kalebu ing~$\Gamma$. Cathetan koreksi sumber OLPL-018:
  sumber Inggris nyebut formula kesimpulan minangka formula kang
  dikenani aturan benar-negasi, sanadyan asumsi pangganti lan aturan
  kasebut mbutuhake formula benar-negasi; kesimpulan iku mapan sawise
  kabeh asumsi, saengga aplikasi lan nomer larik kang dibalekake ing
  kene ditemtokake kanthi tunggal.
\end{proof}

\begin{prop}\ollabel{prop:provability-exhaustive}
  Yen $\Gamma \cup \{!A\}$ lan $\Gamma \cup \{\lnot !A\}$ kalorone
  ora konsisten, mula $\Gamma$ ora konsisten.
\end{prop}

\begin{proof}
  Yen ana $!B_1$, \dots, $!B_n \in \Gamma$ lan $!C_1$, \dots,
  $!C_m \in \Gamma$ saengga
  \begin{align*}
    \{\sFmla{\True}{!A}, &
    \sFmla{\True}{!B_1}, \dots, \sFmla{\True}{!B_n}\} \text{ lan}\\
    \{\sFmla{\True}{\lnot !A}, &
    \sFmla{\True}{!C_1}, \dots, \sFmla{\True}{!C_m}\}
  \end{align*}
  kalorone nduweni !!{tableau}s katutup, kita bisa mbangun siji
  !!{tableau} gabungan kang nuduhake yen $\Gamma$ ora konsisten.
  Gunakna $\sFmla{\True}{!B_1}$, \dots,
  $\sFmla{\True}{!B_n}$ bebarengan karo $\sFmla{\True}{!C_1}$,
  \dots, $\sFmla{\True}{!C_m}$ minangka asumsi, banjur terapna aturan
  \Cut{}. Iki ngasilake rong cabang: siji diwiwiti nganggo
  $\sFmla{\True}{!A}$ lan sijine nganggo $\sFmla{\False}{!A}$.
  
  Ing sisih kiwa, tambahana bagean saka !!{tableau} kapisan kang ana
  ing sangisore asumsi-asumsine. Ing kene, saben aplikasi aturan tetep
  bener, amarga kabeh asumsi !!{tableau} kapisan, kalebu
  $\sFmla{\True}{!A}$, wis cumawis. Mula, saben cabang ing sangisore
  $\sFmla{\True}{!A}$ katutup.
  
  Ing sisih tengen, tambahana bagean saka !!{tableau} kapindho kang
  ana ing sangisore asumsine, kanthi mbusak asil saka saben aplikasi
  $\TRule{\True}{\lnot}$ marang $\sFmla{\True}{\lnot !A}$.
  Kesimpulan $\TRule{\True}{\lnot}$ kang ditrapake marang
  $\sFmla{\True}{\lnot !A}$ yaiku $\sFmla{\False}{!A}$, kang tetep
  cumawis amarga dadi kesimpulan aturan \Cut{} ing sisih tengen
  !!{tableau} gabungan.

  Yen sawijining cabang ing tableau kapindho katutup amarga ngemot
  asumsi $\sFmla{\True}{\lnot !A}$ (kang wis ora muncul minangka
  asumsi ing !!{tableau} gabungan) uga $\sFmla{\False}{\lnot !A}$,
  kita bisa ngetrapake $\TRule{\False}{\lnot}$ marang
  $\sFmla{\False}{\lnot !A}$ kanggo entuk $\sFmla{\True}{!A}$. Saiki
  cabang kang cocog ing !!{tableau} gabungan uga katutup, amarga ngemot
  kesimpulan sisih tengen saka aturan \Cut{}, yaiku
  $\sFmla{\False}{!A}$. Yen sawijining cabang ing !!{tableau}
  kapindho katutup amarga alesan liyane, cabang kang cocog ing
  !!{tableau} gabungan uga katutup, amarga saben !!{signed formula}s
  saliyane $\sFmla{\True}{\lnot !A}$ kang muncul ing cabang
  !!{tableau} lawas kapindho uga muncul ing cabang kang cocog ing
  !!{tableau} gabungan.
  \end{proof}

\end{document}
